% !TEX program = xelatex
\documentclass[10pt,a4paper]{article}

% ======================== 基础包 ========================
\usepackage[top=0.4cm,bottom=0.4cm,left=1.3cm,right=1.3cm]{geometry}
\usepackage{ctex}
\usepackage{titlesec}
\usepackage{enumitem}
\usepackage{xcolor}
\usepackage{hyperref}
\usepackage{fontawesome5}
\usepackage{tikz}

% ======================== 颜色定义 ========================
\definecolor{primary}{HTML}{2B4C7E}
\definecolor{darktext}{HTML}{2D2D2D}
\definecolor{graytext}{HTML}{666666}
\definecolor{rulecolor}{HTML}{B0B0B0}
\definecolor{tagbg}{HTML}{ECF0F5}
\definecolor{tagtext}{HTML}{2B4C7E}

% ======================== 超链接 ========================
\hypersetup{
    colorlinks=true,
    linkcolor=primary,
    urlcolor=primary,
    pdfauthor={徐其栋},
    pdftitle={徐其栋 - AI 应用 / 算法工程师}
}

% ======================== 字体设置 ========================
\setCJKmainfont{SimHei}[BoldFont=SimHei]
\setmainfont{Segoe UI}[BoldFont=Segoe UI Bold]

% ======================== 段落和间距 ========================
\setlength{\parindent}{0pt}
\setlength{\parskip}{0pt}
\pagestyle{empty}

% ======================== 标题样式 ========================
\titleformat{\section}
    {\large\bfseries\color{primary}}
    {}
    {0em}
    {}
    [\vspace{-0.5em}{\color{rulecolor}\rule{\textwidth}{0.6pt}}\vspace{-0.2em}]

\titlespacing*{\section}{0pt}{0.35em}{0.15em}

% ======================== 自定义命令 ========================
\newcommand{\projectblock}[3]{%
    {\bfseries\color{darktext}#1}\hfill{\color{graytext}#2}\\%
    {\small\color{graytext}#3}%
    \vspace{0.2em}%
}

% ======================== 文档开始 ========================
\begin{document}

% ======================== 头部 ========================
\begin{center}
    {\fontsize{26}{32}\selectfont\bfseries\color{primary} 徐其栋}

    \vspace{0.15em}

    {\color{graytext}
    \faPhone\hspace{0.3em}17268891140
    \hspace{1.0em}
    \faEnvelope\hspace{0.3em}\href{mailto:xqd922@outlook.com}{xqd922@outlook.com}
    \hspace{1.0em}
    \faGithub\hspace{0.3em}\href{https://github.com/xqd922}{github.com/xqd922}
    \hspace{1.0em}
    \faGlobe\hspace{0.3em}\href{https://1001.pp.ua}{个人主页}}

    \vspace{0.5em}

    {\color{darktext} 求职意向：\textbf{\color{primary}AI 应用工程师 / 算法工程师（初级） / AI 辅助开发}}
\end{center}

\vspace{0.2em}
{\color{rulecolor}\hrule height 0.8pt}

% ======================== 教育背景 ========================
\section{\faGraduationCap\hspace{0.4em}教育背景}

{\bfseries 滇西科技师范学院}\hfill{\bfseries 物联网工程\hspace{0.3em}|\hspace{0.3em}本科}

{\color{graytext} 2022.09 -- 2026.06}

\vspace{0.2em}

{\color{darktext}
\textbf{核心课程：}Python 程序设计、人工智能、图像处理、数据库应用、机器学习基础、数据结构、计算机网络
}

% ======================== 专业技能 ========================
\section{\faCogs\hspace{0.4em}专业技能}

{\bfseries\color{primary}编程语言与数据处理：}\hspace{0.3em}{\color{darktext}Python、NumPy、OpenCV、Matplotlib、C 语言、SQL}

\vspace{2pt}

{\bfseries\color{primary}人工智能与图像处理：}\hspace{0.3em}{\color{darktext}中值滤波、直方图均衡化、拉普拉斯锐化、二维卷积、图像增强流水线、人工智能课程基础}

\vspace{2pt}

{\bfseries\color{primary}LLM 与 AI 工程：}\hspace{0.3em}{\color{darktext}ChatGPT、Claude、DeepSeek、OpenAI / Claude API 调用、Prompt Engineering、Claude Code Skills、AI Pair Programming}

\vspace{2pt}

{\bfseries\color{primary}前端与工程：}\hspace{0.3em}{\color{darktext}Next.js 15、React 19、TypeScript、MDX、Tailwind CSS、Vercel 部署、Git、Linux、MySQL}

% ======================== 项目经历 ========================
\section{\faProjectDiagram\hspace{0.4em}项目经历}

% --- 机器狗（AI 交互重写） ---
\projectblock{具备语音识别交互的四足仿生机器狗}{2025.03 -- 2025.06}{技术栈：ESP32 语音识别\hspace{0.3em}|\hspace{0.3em}STM32\hspace{0.3em}|\hspace{0.3em}运动学逆解\hspace{0.3em}|\hspace{0.3em}多模态交互\hspace{0.3em}|\hspace{0.3em}蓝牙}

\begin{itemize}[leftmargin=1.5em, itemsep=2pt, parsep=0pt, topsep=2pt]
    \item \textbf{AI 语音交互模块}：基于 ESP32 集成语音识别模块，将"前进 / 后退 / 左转 / 挥手"等自然语言指令解析为结构化控制命令，实现语音到动作的端到端映射。
    \item \textbf{运动智能算法}：基于余弦定理与三角函数完成腿部运动学逆解，由足尖目标坐标 $(x,y)$ 反推关节角度，配合步进插值算法实现动作平滑过渡，共计 9 类仿生动作。
    \item \textbf{多模态交互}：打通语音识别 + 蓝牙网页遥控 + 状态灯反馈三通道交互链路，具备 AI 感知、指令理解、运动执行的完整闭环能力。
\end{itemize}

\vspace{0.15em}

% --- 个人博客（AI 辅助开发重写） ---
\projectblock{AI 辅助开发的个人技术博客}{2024.11 -- 至今}{技术栈：Next.js 15\hspace{0.3em}|\hspace{0.3em}React 19\hspace{0.3em}|\hspace{0.3em}TypeScript\hspace{0.3em}|\hspace{0.3em}MDX\hspace{0.3em}|\hspace{0.3em}Tailwind CSS\hspace{0.3em}|\hspace{0.3em}Vercel\hspace{0.3em}|\hspace{0.3em}Claude / GPT API}

\begin{itemize}[leftmargin=1.5em, itemsep=2pt, parsep=0pt, topsep=2pt]
    \item \textbf{AI 协同开发}：以 Claude Code / ChatGPT 作为结对编程助手，通过精细化 Prompt 设计与 Claude Skills 复用模式，完成 UI 组件、页面交互、动画与样式的高效生成与迭代。
    \item \textbf{LLM API 集成实践}：实际使用 OpenAI / Claude API 完成文本生成、内容摘要、代码辅助等调用场景，熟悉 API 调用参数（temperature、max\_tokens、system prompt）与流式响应处理。
    \item \textbf{工程化上线}：基于 Next.js 15 + React 19 + TypeScript + MDX 构建静态博客，Tailwind CSS 样式、Framer Motion 动画，Vercel 持续部署，源码托管于 \href{https://github.com/xqd922/blog}{github.com/xqd922/blog}，访问域名 \href{https://1001.pp.ua}{1001.pp.ua}。
\end{itemize}

\vspace{0.15em}

% --- 数字图像处理综合实践项目（真实小组课程项目） ---
\projectblock{数字图像处理综合实践（车牌图像增强小组项目）}{2025.05}{技术栈：Python\hspace{0.3em}|\hspace{0.3em}OpenCV\hspace{0.3em}|\hspace{0.3em}NumPy\hspace{0.3em}|\hspace{0.3em}中值滤波\hspace{0.3em}|\hspace{0.3em}直方图均衡化\hspace{0.3em}|\hspace{0.3em}拉普拉斯锐化}

\begin{itemize}[leftmargin=1.5em, itemsep=2pt, parsep=0pt, topsep=2pt]
    \item \textbf{协作与分工}：5 人小组完成车牌图像修复流水线（椒盐噪声去除 $\to$ 暗处细节增强 $\to$ 模糊锐化），独立负责\textbf{图像模糊改善}模块的算法原理、代码实现与效果评估。
    \item \textbf{拉普拉斯锐化算法实现}：基于 OpenCV 的 \texttt{cv2.filter2D} 实现 3$\times$3 拉普拉斯锐化核（中心系数 $+5$）的二维卷积，强化图像边缘与高频细节，显著提升车牌字符可识别度。
    \item \textbf{流水线协同}：与中值滤波去噪、直方图均衡化提亮模块级联对接，形成多算法图像增强流水线；熟悉 OpenCV 图像 I/O、灰度转换、核卷积与结果对比分析全流程。
\end{itemize}

\vspace{0.15em}

% ======================== 个人评价 ========================
\section{\faUser\hspace{0.4em}个人评价}

2026 届本科生，求职方向 AI 应用 / 初级算法。熟悉 Python 与 OpenCV 图像处理，独立完成过拉普拉斯锐化等经典图像增强算法落地；实际使用过 OpenAI / Claude API，熟练 Prompt 工程、Claude Skills 与 AI Pair Programming；具备 Next.js + TypeScript + MDX 前端工程与 MySQL 数据库能力，可快速进入 AI 应用研发工作。

\vspace{0.15em}

% ======================== 底部标签 ========================
\begin{center}
\tikz{\node[fill=tagbg, rounded corners=3pt, inner sep=5pt, text=tagtext, font=\small\bfseries] {Python};}
\hspace{0.2em}
\tikz{\node[fill=tagbg, rounded corners=3pt, inner sep=5pt, text=tagtext, font=\small\bfseries] {OpenCV};}
\hspace{0.2em}
\tikz{\node[fill=tagbg, rounded corners=3pt, inner sep=5pt, text=tagtext, font=\small\bfseries] {图像处理};}
\hspace{0.2em}
\tikz{\node[fill=tagbg, rounded corners=3pt, inner sep=5pt, text=tagtext, font=\small\bfseries] {Claude / GPT API};}
\hspace{0.2em}
\tikz{\node[fill=tagbg, rounded corners=3pt, inner sep=5pt, text=tagtext, font=\small\bfseries] {Prompt 工程};}
\hspace{0.2em}
\tikz{\node[fill=tagbg, rounded corners=3pt, inner sep=5pt, text=tagtext, font=\small\bfseries] {Next.js / TS};}
\hspace{0.2em}
\tikz{\node[fill=tagbg, rounded corners=3pt, inner sep=5pt, text=tagtext, font=\small\bfseries] {Git / Linux};}
\end{center}

\end{document}
